% SOURCE_AUTHORITY: sga5_fr_workpass.tex, lines 7268--7291 (LF-normalized unit).
% SOURCE_SHA256: 791F4EFFC5E02832D5D77ED03518C8156D6F07E4C8238B03545DB93D883FBB28
\begin{statement}{Definición 4.1.2}
Sea $X$ un sistema proyectivo indexado por $\ZZ$. Se dice que $X$ es esencialmente constante si existe un entero $p$ tal que, para todo par $(r,s)$ de enteros mayores o iguales que $p$ y que verifican la desigualdad $s\ge r$, el morfismo de transición
\[
X_s\longrightarrow X_r
\]
sea un isomorfismo. Es claro que basta exigirlo cuando $r=p$.
\end{statement}

\begin{statement}{Proposición 4.1.3 (Propiedad (MLAR) y graduados)}
Las afirmaciones siguientes para un objeto $X$ de la categoría $P$ son equivalentes:
\begin{enumerate}[label=(\roman*)]
\item $X$ verifica (MLAR).
\item El sistema proyectivo con valores en la categoría de los objetos de $\cC$ graduados por $\NN$
\[
r\longmapsto (\gr^\alpha(X_{\alpha+r}))_{\alpha\in\NN},
\]
donde los morfismos de transición son evidentes, es esencialmente constante.
\item El sistema proyectivo con valores en la categoría de los objetos de $\cC$ graduados por $\NN$
\[
n\longmapsto (\gr^p(X_n))_{p\in\NN},
\]
donde los morfismos de transición son evidentes, verifica (MLAR).
\end{enumerate}
\end{statement}
